\section{Model-Based RL 补充：合成数据生成}

本讲的前半部分先把上一讲尚未讲完的 model-based RL 收尾。Chelsea Finn 的切入点很明确：如果我们已经学到了一个能够预测未来状态的模型，那么它除了拿来做 planning，还能不能用来\textbf{制造更多训练数据}，从而反过来训练一个更强、更便宜的策略？

\begin{figure}[H]
\centering
\includegraphics[width=0.9\textwidth]{slides-images/slide-04.jpg}
\caption{Planning with learned models：规划、收集数据、更新模型、周期性重规划}
\end{figure}

\subsection{从 planning 走向 policy optimization}

上一讲讲的是 ``给定当前状态，用 learned model 在测试时搜索未来动作序列''。这一讲更进一步：如果测试时反复搜索太贵，能否把 model 提供的预测能力转成策略学习阶段的数据增广工具？

\begin{importantbox}{这一节真正要解决的问题}
planning 与 policy optimization 的分工不同：
\begin{itemize}
    \item planning 的优点是灵活，目标换了仍能重算。
    \item planning 的缺点是 test-time compute 高，而且长 horizon 时模型误差累积明显。
    \item 直接学策略的优点是部署便宜。
    \item 直接学策略的缺点是如果训练数据不够，策略很难学稳。
\end{itemize}
因此，最自然的问题就是：\textbf{如何用 learned model 为策略制造更多有效数据，同时避免把模型误差放大到无法控制。}
\end{importantbox}

\subsection{Dyna 风格的数据增强}

\begin{figure}[H]
\centering
\includegraphics[width=0.9\textwidth]{slides-images/slide-05.jpg}
\caption{Model-based policy optimization 的核心选择：从哪里开始 rollout，roll 多长}
\end{figure}

课程把这个问题讲得非常具体。假设我们已经在真实环境中收集到若干轨迹，那么模型生成数据时至少有三种策略：
\begin{enumerate}
    \item 从初始状态开始 rollout 很长的轨迹。
    \item 只从初始状态开始做短 rollout。
    \item 从数据集中\textbf{所有已出现过的状态}出发做短 rollout。
\end{enumerate}

第三种策略是本讲推荐的核心做法，因为它兼顾了覆盖面与误差控制。

\begin{figure}[H]
\centering
\includegraphics[width=0.9\textwidth]{slides-images/slide-06.jpg}
\caption{更稳妥的折中：从数据里的所有状态出发，生成局部 synthetic rollout}
\end{figure}

\begin{knowledgebox}{为什么 ``短 rollout + 所有状态'' 是最稳妥的}
\begin{enumerate}
    \item 如果只从初始状态生成长轨迹，模型误差会在长链条上不断放大。
    \item 如果只从少数起点生成短轨迹，后续状态覆盖又会严重不足。
    \item 从数据中的所有状态出发做短 rollout，本质上是在真实分布附近做局部外推，最能兼顾稳定性与覆盖度。
\end{enumerate}
\end{knowledgebox}

\begin{figure}[H]
\centering
\includegraphics[width=0.9\textwidth]{slides-images/slide-07.jpg}
\caption{完整算法：真实数据更新模型，模型短 rollout 生成额外训练样本}
\end{figure}

\subsection{什么时候值得学世界模型}

model-based RL 的诱惑很强，因为它看起来像 ``会想象未来''。但 Chelsea Finn 的判断标准很克制：\textbf{关键不在于模型是否听起来高级，而在于它是否真的比策略更容易学。}

\begin{figure}[H]
\centering
\includegraphics[width=0.9\textwidth]{slides-images/slide-08.jpg}
\caption{Model-based RL 的优点与代价：数据效率、任务无关性与训练复杂度的平衡}
\end{figure}

讲者给出的判断原则可以概括成一张简表：

\begin{table}[H]
\centering
\begin{tabular}{p{0.2\textwidth} p{0.28\textwidth} p{0.28\textwidth}}
\toprule
场景特征 & 更可能有利于 model-based & 风险点 \\
\midrule
真实交互昂贵 & 可以靠 synthetic rollout 提高数据效率 & 模型偏差被策略吸收 \\
动力学相对规则 & 世界模型较容易学出可用近似 & 容易低估长 horizon 误差 \\
奖励经常变化 & 同一模型可服务多个下游目标 & 不代表模型本身优化了任务表现 \\
环境高度随机或部分可观测 & 未必适合 & 可能比直接学策略更难 \\
\bottomrule
\end{tabular}
\caption{何时优先考虑 model-based RL}
\end{table}

\begin{warningbox}{不要把 model-based 当成默认更强的路线}
世界模型需要单独训练、单独调参，还会引入新的误差来源。若环境动力学本身比策略更难学，或者 rollout 稍长就会失真，那么学模型反而可能让整个系统更脆弱。
\end{warningbox}

\subsection{世界模型不只有一种形式}

课程最后提醒：传统 forward dynamics $p(s_{t+1}\mid s_t, a_t)$ 只是世界模型的一种。inverse model、多步 inverse model、只预测未来观测的模型、甚至 transition tuple 的联合分布都可能在特定任务里更有用。

\begin{figure}[H]
\centering
\includegraphics[width=0.9\textwidth]{slides-images/slide-09.jpg}
\caption{世界模型的多种形态：forward/inverse/future prediction 各有用途}
\end{figure}

\subsection{本章小结}

model-based RL 在这一讲里被重新定位为一种\textbf{数据增强工具}：关键不是做最完美的世界模拟，而是用足够准确的局部模型在真实数据附近扩充训练样本，并始终把模型偏差控制在可接受范围内。

